
\newpage
\addcontentsline{toc}{chapter}{\MakeUppercase{Appendix III: Sample Code}}
\thispagestyle{empty}
%
\begin{spacing}{1.5}
%
\section*{Appendix III: Sample Code}
%%
%% Here enter abstract. No figures, sketches, tables shall be here.
%%

here you may add the sample code of your project

\begin{enumerate}
\item Merge Sort Function in C
\begin{verbatim}
// ignore_for_file: prefer_const_constructors

import 'package:cloud_firestore/cloud_firestore.dart';
import 'package:firebase_auth/firebase_auth.dart';
import 'package:flutter/material.dart';
import 'package:quick_bite/pages/ownerhome.dart';
class AddBranch extends StatelessWidget {
  const AddBranch({super.key});

  signup(String name,String branch,String company,String phone,String email,String password,String confirm,BuildContext context) async {
    if(passwordConfirmed(password,confirm)){
      await FirebaseAuth.instance.createUserWithEmailAndPassword(
        email: email.trim() ,
       password: password.trim());
        final user= FirebaseAuth.instance.currentUser;
        final uid = user?.uid;
       addUserDetails(
        name.trim(),
        branch.trim(),
        company.trim(),
        phone.trim(),
        email.trim(),
        uid,
        context,
       );
    }
  }

  addUserDetails(String name,String branch,String company,String phone,String email,String? uid, context) async{
    Map<String, dynamic> UserData ={
      'name': name,
      'branch': branch,
      'company': company,
      'phoneno': phone,
      'role': 'manager',
      'staffcount':'0'
    };
   
    FirebaseFirestore.instance.collection('users').doc(uid).set(UserData);
  }

  bool passwordConfirmed(String password,String confirm) { 
    if(password.trim() == confirm.trim()){
      return true;
    }else{
      return false;
    }
  }

  Future<String> detailscompany() async {
  try {
    final compuser = FirebaseAuth.instance.currentUser;
    final ownuid = compuser?.uid;
    DocumentSnapshot snapshot =
        await FirebaseFirestore.instance.collection("users").doc(ownuid).get();
    if (snapshot.exists && snapshot.data() != null) {
      Map<String, dynamic> userData = snapshot.data() as Map<String, dynamic>;
      String company = userData['company'];
      return company;
    } else {
      throw Exception('User not found');
    }
  } catch (e) {
    print('Error fetching company details: $e');
    throw Exception('Failed to fetch company details');
  }
}


  @override
  Widget build(BuildContext context) {
    final user = FirebaseAuth.instance.currentUser;
    final uid = user?.uid;
    TextEditingController name=TextEditingController();
    TextEditingController branch=TextEditingController();
    TextEditingController phone=TextEditingController();
    TextEditingController email=TextEditingController();
    TextEditingController password=TextEditingController();
    TextEditingController confirm=TextEditingController();

    return StreamBuilder(
    stream: FirebaseFirestore.instance.collection('users').doc(uid).snapshots(),
    builder: (context, snapshot) {
      if (snapshot.connectionState == ConnectionState.waiting) {
        return CircularProgressIndicator(); // Show loading indicator while fetching company details
      } else if (snapshot.hasError) {
        return Text('Error: ${snapshot.error}'); // Show error if fetching company details fails
      } else {
        final data = snapshot.data!.data();
        final company = data!['company'];

    return Scaffold(
      backgroundColor: Colors.grey[300],
      appBar: AppBar(
        title: Center(
          child: Text("Enter Details",style: TextStyle(
          fontSize: 30.0,
          color: Colors.black,
          fontWeight: FontWeight.bold
        ),),
        ),
      backgroundColor: Colors.white,
      automaticallyImplyLeading: true,
      ),
      body: Center(
        child: ListView(
          children: [
            SizedBox(
              height: 20,
            ),
            Padding(padding: EdgeInsets.all(15),
              child:  TextField(
                controller: branch,
                decoration: InputDecoration(
                  border: OutlineInputBorder(borderRadius: BorderRadius.circular(12)),
                  labelText: "enter branch name:"
                ),
              ),
             ),
              SizedBox(
              height: 5,
            ),
            Padding(padding: EdgeInsets.all(15),
              child:  TextField(
                controller: name,
                decoration: InputDecoration(
                  border: OutlineInputBorder(borderRadius: BorderRadius.circular(12)),
                  labelText: "enter branch manager name:"
                ),
              ),
             ),
              SizedBox(
              height: 5,
            ),
            
            Padding(padding: EdgeInsets.all(15),
              child:  TextField(
                controller: phone,
                decoration: InputDecoration(
                  border: OutlineInputBorder(borderRadius: BorderRadius.circular(12)),
                  labelText: "enter manager phone number:"
                ),
              ),
             ),
             SizedBox(
              height: 5,
            ),
            Padding(padding: EdgeInsets.all(15),
              child:  TextField(
                controller: email,
                decoration: InputDecoration(
                  border: OutlineInputBorder(borderRadius: BorderRadius.circular(12)),
                  labelText: "enter manager email:"
                ),
              ),
             ),
             SizedBox(
              height: 5,
            ),
            Padding(padding: EdgeInsets.all(15),
              child:  TextField(
                controller: password,
                decoration: InputDecoration(
                  border: OutlineInputBorder(borderRadius: BorderRadius.circular(12)),
                  labelText: "enter password:"
                ),
              ),
             ),
             SizedBox(
              height: 5,
            ),
            Padding(padding: EdgeInsets.all(15),
              child:  TextField(
                controller: confirm,
                decoration: InputDecoration(
                  border: OutlineInputBorder(borderRadius: BorderRadius.circular(12)),
                  labelText: "confirm password:"
                ),
              ),
             ),
             GestureDetector(
                onTap: () {
                signup(name.text,branch.text,company,phone.text,email.text,password.text,confirm.text,context);
                Navigator.push(context, MaterialPageRoute(builder: ((context) => OwnerHome())));
                },
                child: Padding(
                  padding: const EdgeInsets.all(15.0),
                  child: Container(
                    decoration: BoxDecoration( 
                      color: Colors.grey[700],
                      borderRadius: BorderRadius.circular(12),
                    ),
                    padding: EdgeInsets.all(25),
                    child: Center(
                      child: Text( 
                        'CREATE',
                        style: TextStyle(
                          color: Colors.white,
                          fontWeight: FontWeight.bold,
                          fontSize: 16,
                        ),
                      ),
                    ),
                  ),
                ),
              )
          ],
        ),
      ),
    );
  }
}
    );}}
\end{verbatim}
\item Merge Function in C

\begin{verbatim}
void merge(int arr[], int l, int m, int r)
{
int i, j, k;
int n1 = m - l + 1;
int n2 = r - m;
// Create temp arrays
int L[n1], R[n2];
// Copy data to temp array
for (i = 0; i < n1; i++)
L[i] = arr[l + i];
for (j = 0; j < n2; j++)
R[j] = arr[m + 1+ j];
// Merge the temp arrays
i = 0;
j = 0;
k = l;
while (i < n1 && j < n2)
{
if (L[i] <= R[j])
{
arr[k] = L[i];
i++;
}
else
{
arr[k] = R[j];
j++;
}
k++;
}
// Copy the remaining elements of L[]
while (i < n1)
{
arr[k] = L[i];
i++;
k++;
}
// Copy the remaining elements of R[]
while (j < n2)
{
arr[k] = R[j];
j++;
k++;
}
}
\end{verbatim}
\end{enumerate}
%%
\end{spacing}
